\section{Discussion}
\label{sec:discussion}

\mra{} treats underwater gate racing as a task in ordered-sequence autonomy
rather than in speed alone, and its design follows from that choice. A vehicle
must commit to the correct target, reacquire a new target after each passage and
compose those transitions over a long course, so the object of evaluation is the
chain of transitions rather than any single passage. The benchmark therefore
specifies ordered circuits of twelve, seventeen and twenty-two gates instead of
an isolated gate-passing exercise, and charges the referee with validating the
order of crossings rather than their count alone.

\subsection{What the Evaluation Contract Exposes}

The central value of \mra{} is the conjunction of a fixed information boundary,
an ordered referee, configurable race management and machine-readable scoring.
The boundary makes controller comparisons like-for-like: rule-based and learned
controllers operate on the same participant-level information. The referee makes
local progress auditable: controller estimates can be compared with privileged
adjudication without allowing the latter into the control loop. The resulting
residual disagreements are useful evidence, not hidden implementation details.

The configurable race manager makes this separation operational. Track
geometry, participants, environmental conditions, starts, timing and penalties
can be changed through scenario configuration while the participant interface
and referee logic remain fixed. This permits the benchmark to vary the task and
the operating condition independently of the autonomy implementation, which is
important when interpreting differences between controllers.

The referee also makes failure interpretable. A missed next gate is recorded as
a sequence failure rather than being blurred into a low aggregate gate count,
and safety costs remain visible alongside progress. Completion, progress and
safety events are reported jointly for the same reason: read together they make
the effect of environmental disturbance visible beyond completion rate alone.
The medium-current condition illustrates the point, since under it both
controllers completed fewer runs and recorded collisions where the clean
condition recorded none.

\subsection{Reference Controllers and Fleets}

The reference-controller comparison places two controllers that differ in a
documented respect on the same circuits under the same protocol. Neither
recorded the better outcome everywhere. On the long mixed course, suspending
late visual corrections completed all five seeds where continuous correction
completed two; on the repeated vertical transitions a single run separates the
two in the opposite direction. These are outcomes recorded over a small number
of seeds, not orderings the benchmark establishes.

The fleet experiments extend the same logic from individual to team autonomy.
Vehicles that are individually competent can still collide when their
interaction is not scored. With a distributed leader--follower policy operating
within the legal observation boundary, no such events were recorded at either
release gap. The two variants are not charged alike, however: LF(1) is charged
nothing in either configuration, whereas LF(2) incurs one stuck-timer charge per
run at simultaneous release and none at the $8$\,s gap. Coordination and
whatever it costs are therefore both carried by the official penalized team
score rather than by a separate engineering trace.

\subsection{Learning as a Test of Interface Generality}

The learning experiment exercises the generality of the interface rather than of
any learned policy. A recurrent PPO controller can be integrated within the same
participant-level information boundary, body-frame command interface and
independent referee procedure as the reference controllers, and can be assessed
on complete ordered circuits under the reported current-free protocol. The policy
was trained on the same three official circuit families on which it is then
evaluated, so these episodes speak to integration and in-distribution execution
rather than to transfer to unseen geometry. The interface constrains the legal information and
action contract and does not prescribe the internals of a controller: the present
work instantiates it with rule-based and learned controllers, and any other
autonomy strategy that respects the same contract can be evaluated through it
without changes to the race-management or referee layers.
